%% sections/02_model.tex
\section{The Delay Differential Equation Model}
\label{sec:model}

The model, originally proposed by Villasana and Radunskaya~\cite{villasana2003} and extended in~\cite{villasana2004}, tracks four coupled populations: $T_I(t)$ (tumor cells in interphase), $T_M(t)$ (tumor cells in mitosis), $I(t)$ (immune cells), and $u(t)$ (drug concentration).

\subsection{Assumptions}

All mathematical models of biological systems are simplifications. The key assumptions here are~\cite{villasana2003}: (1) mitosis is instantaneous relative to the overall cell cycle duration; (2) the drug specifically targets and arrests tumor cells during mitosis; (3) the drug is cytotoxic to both tumor and immune cells, reflecting the bone marrow toxicity common to many chemotherapy agents; (4) the quiescent $G_0$ phase is excluded from the base model, despite quiescent cells being estimated to comprise roughly 20\% of total tumor mass and are largely resistant to chemotherapy. This exclusion is addressed in Section~\ref{sec:extension}; (5) drug concentration decays exponentially, corresponding to a single-compartment pharmacokinetic model.

\subsection{System of Equations}

\begin{align}
  T_I' &= 2a_4T_M - (c_1I + d_2)T_I - a_1T_I(t-\tau) \label{eq:TI} \\
  T_M' &= a_1T_I(t-\tau) - (d_3 + a_4 + c_3I)T_M \notag \\
       &\quad - k_1(1-e^{-k_2u})T_M \label{eq:TM} \\
  I'   &= k + \frac{\rho I(T_I+T_M)^n}{\alpha+(T_I+T_M)^n} \notag \\
       &\quad - (c_2T_I + c_4T_M + d_1 + k_3(1-e^{-k_2u}))I \label{eq:I} \\
  u'   &= -\gamma u \label{eq:u}
\end{align}

The delay $\tau$ in Equations~\ref{eq:TI} and~\ref{eq:TM} represents the time spent by tumor cells in interphase before entering mitosis. The term $a_1T_I(t-\tau)$ captures cells that entered interphase $\tau$ time units ago and are now transitioning to division. This is the key biological realism that \gls{ode} models cannot represent.

In Equation~\ref{eq:I}, the Michaelis-Menten term models immune stimulation by the tumor: as tumor burden increases, immune response grows up to a saturation ceiling $\alpha$. The drug term $k_3(1-e^{-k_2u})$ captures immune cell toxicity from chemotherapy.

\subsection{Biological Interpretation of Terms}

\begin{align*}
  T_I' &= \underbrace{2a_4T_M}_{\substack{\text{cells entering}\\\text{interphase}\\\text{after mitosis}}}
         -\underbrace{(c_1I + d_2)T_I}_{\substack{\text{combat \&}\\\text{natural deaths}}}
         -\underbrace{a_1T_I(t-\tau)}_{\substack{\text{cells entering}\\\text{mitosis (delayed)}}} \\[6pt]
  T_M' &= \underbrace{a_1T_I(t-\tau)}_{\substack{\text{cells entering}\\\text{mitosis}}}
         -\underbrace{(d_3+a_4)T_M}_{\substack{\text{natural death}\\\text{\& exit to }T_I}}
         -\underbrace{c_3T_MI}_{\text{combat deaths}}
         -\underbrace{k_1(1-e^{-k_2u})T_M}_{\text{drug effect}}
\end{align*}

\subsection{Parameters}

\begin{table}[H]
  \caption{Parameter values from~\cite{villasana2004} (non-dimensional). We use
    $\gamma := \lambda_2 = 0.85$ following the single-drug-term base
    model~\cite{villasana2003}.}
  \label{tab:params}
  \small
  \begin{tabular}{llll}
    \toprule
    $\tau = 0.92$ & $a_1 = 0.98$   & $a_4 = 0.8$     & $d_2 = 0.11$  \\
    $d_3 = 0.4$   & $d_1 = 0.29$   & $c_1 = 0.9$     & $c_3 = 0.9$   \\
    $c_2 = 0.085$ & $c_4 = 0.085$  & $\rho = 0.1$    & $\alpha = 0.2$\\
    $k = 0.036$   & $n = 3$        & $k_1 = 0.47$    & $k_2 = 0.57$  \\
    $k_3 = 0.49$  & $\gamma = 0.85$&                 &               \\
    \bottomrule
  \end{tabular}
\end{table}
